\section{Problem and Requirements}
\label{sec:motivation}

\subsection{The detection task}
A \sysname{} agent reports metrics on a fixed cadence: CPU and memory for a host,
octet and error counters for an interface, latency and availability for a service.
Each distinct stream is one time series. The task is to flag, in near real time,
when a series departs from its own recent behavior. We target point and level
anomalies, where a sample is far from the recent baseline, scored per series
against that series' own history. Detection is independent across series: there is
no cross-series correlation at this layer. This engine is the fast, high-recall
tier of a two-tier design: it proposes candidate anomalies the moment a series
jumps. A short window cannot tell a real problem from a normal busy Tuesday
morning, when the recent window is full of overnight-quiet samples and the morning
ramp reads as a breach. A complementary central tier answers that seasonal question
by scoring each series against a profile of its own history at the same time of
week, and disposes of the edge's candidates accordingly; we describe that tier in
\cref{sec:seasonal}.

\begin{figure}[t]
  \centering
  \fitw{%
  \begin{tikzpicture}[font=\small]
    \node[box, text width=32mm] (edge) at (0,1.7)
      {\textbf{Edge} (per agent)\\ fast spike, \emph{high recall}\\ rolling $z$-score};
    \node[box, text width=32mm] (cen) at (0,-1.7)
      {\textbf{Central} (per tenant)\\ seasonal, \emph{precision}\\ hour-of-week profile};
    \node[box, text width=20mm] (join) at (5.6,0) {verdict join\\(series, time)};
    \node[box, text width=33mm, align=left] (out) at (9.9,0)
      {suppress seasonal-expected;\\ escalate spike \& off-baseline;\\ surface seasonal-only};
    \draw[ar] (edge.east) -- node[midway,above,font=\scriptsize]{candidates} (join.north west);
    \draw[ar] (cen.east) -- node[midway,below,font=\scriptsize]{context} (join.south west);
    \draw[ar] (join.east) -- (out.west);
  \end{tikzpicture}}
  \caption{The two tiers compose as recall and precision: the edge proposes spike
  candidates, and the central seasonal tier disposes by joining on series and time
  window. The central tier is outside this paper's scope.}
  \label{fig:tiers}
\end{figure}

\subsection{Workload characteristics}
Two properties of the workload shape the design. First, series are numerous and
mostly independent, so per-series cost is what matters, and a fixed per-series
budget multiplied by the series count must stay small. Second, metrics come in two
shapes. Gauges (CPU percent, queue depth) can be scored directly. Cumulative
monotonic counters (interface octets, packet counts) cannot: their raw value only
ever climbs, so a $z$-score of the raw counter is meaningless. They have to be
converted to a per-second rate first. We also exclude per-process metrics from
detection, because their cardinality is unbounded and dominated by short-lived
processes that produce no useful baseline.

\subsection{Requirements}
These lead to six requirements that the rest of the paper returns to.
\begin{itemize}
  \item \textbf{R1, per-series independence.} One detector state per series, with
        no coupling between series.
  \item \textbf{R2, stateless reasoning over an explicit context.} A pass takes a
        context and a sample and returns a verdict plus the next state, so state is
        easy to snapshot and restore.
  \item \textbf{R3, bounded resources.} A fixed memory cost per series and a hard
        cap on the number of series, so one busy host cannot grow the engine
        without limit.
  \item \textbf{R4, counter awareness.} Cumulative counters are rate-normalized
        before scoring, with correct handling of resets, wraps, and gaps.
  \item \textbf{R5, cheap per evaluation.} Each sample is cheap enough to score as
        it arrives, so the engine keeps up in real time on a small CPU budget.
  \item \textbf{R6, actionable verdicts.} A verdict is a structured finding, keyed
        so it can be correlated and alerted on downstream.
\end{itemize}
